% =============================================================================
\section{Binding code}
\label{sec:code}
% =============================================================================
Boost.Python allows users to expose C++ classes and functions to Python using nothing more than a C++ compiler. It exploits metaprogramming techniques to implement a rich set of features and high-level user-friendly interface. Boost.Python extracts as much as information as possible from the source code to be wrapped. This approach is referred to as \emph{user guided wrapping}. Additional information that cannot be automatically deduced is explicitly supplied by the user. The interface specification is written in the same full-featured C++ language as the code being exposed, which has the potential of generating efficient binding. In this section we describe some of the techniques we use to ensure that.

% -----------------------------------------------------------------------------
\subsection{Following Generic Concepts}
\label{sec:code:concepts}
% -----------------------------------------------------------------------------
As mentioned in the introduction, \cgal{} rigorously adheres to the generic programming paradigm. As a consequence, most components of \cgal{} are either class or function templates. Many of the parameters of these templates are described in terms of models of concepts. When a class or a function template is instantiated, each one of its template parameters is substituted with a model of one or more concepts associated with the template parameter. These concepts are organized in clusters that describe the refinement relations among the concepts within a cluster. Close to~750 concepts can be identified in \cgal{} at the time this article is written. Most hierarchies are small. Few hierarchies, such as the hierarchy of concepts of the geometry traits of the \iiDArrangementsPackage{} package is quite large with intricate refinement relations. We have introduced a tight coupling between concepts and binding generations. In particular, we have introduced a template function that generates bindings for each concept. The function generates the bindings for all the type and function members listed in the concept. We call this function for every model of the concept. This systematic approach guaranties that all the documented functions and types, but nothing else, are exposed. Consider for example the cluster hierarchy of geometry traits concepts depicted in Figure~\ref{fig:atc:central}. The template functions \cppLstinline{export_AosBasicTraits\_2()}, \cppLstinline{export_AosXMonotoneTraits\_2()}, and \cppLstinline{export_AosTraits\_2()} accept a binding class object for a particular geometry traits model and populate it with all attributes that correspond to the requirements of the concepts \cppLstinline{AosBasicTraits\_2}, \cppLstinline{AosXMonotoneTraits\_2}, and \cppLstinline{AosTraits\_2}, respectively. If a concept \cppLstinline{B} refines a concept \cppLstinline{A}, then the function \cppLstinline{export_B()} calls the function \cppLstinline{export_A()}. The bindings of different traits models are generated in separate compilation units. As nodes in the concept refinement graphs may have more than a single parent, we ensure that a function that corresponds to a concept is not invoked more than once in a given compilation unit; see Lines~3 and 4 in the definition of the function \cppLstinline{export_AosBasicTraits\_2()} below.

\begin{lstlisting}[style=cppNumberedStyle,basicstyle=\scriptsize]
template <typename GeometryTraits_2, typename ClassObject, typename Concepts>
void export_AosTraits_2(ClassObject c, Concepts& concepts) {
  namespace bp = boost::python;

  static bool exported = false;
  if (exported) return;

  export_AosXMonotoneTraits_2<GeometryTraits_2>(c, concepts);

  bp::scope traits_scope(c);
  auto& classes = concepts.m_basic_traits_classes;

  // Attribute added to the traits scope follow
  classes.m_point_2 = bp::class_<Point_2>("Point_2")
    .def(bp::init<>())
    .def(bp::init<const Point_2&>());
}
\end{lstlisting}

Typically, a concept in our case requires the provision of nested types that are themselves models of other concepts. For example, the concept \cppLstinline{AosBasicTraits\_2} requires the provision of the nested type \cppLstinline{Point_2} used to represent a two-dimensional point. The type \cppLstinline{Point_2} must be \cppLstinline{DefaultConstructible} and \cppLstinline{CopyConstructible}. Let \cppLstinline{p} be an object of type  \cppLstinline{Point_2} required by the \cppLstinline{AosBasicTraits\_2} concept. Most traits models also support the calls \cppLstinline{p.x()} and \cppLstinline{p.y()}; they return the $x$ and $y$ exact Cartesian coordinates of the points $p$, respectively. However, the \cppLstinline{Point_2} type nested in the traits model \cppLstinline{Arr_Bezier_curve_traits\_2} does not maintain an exact representation of the coordinates. Thus, it instead supports the call \cppLstinline{p.approximate()}, which returns a pair of floating point numbers that only approximate the real coordinates. The function \cppLstinline{export_AosBasicTraits\_2()} constructs a class object for the \cppLstinline{Point\_2} type, but adds Python attributes only for the default and copy constructors of this type; see Lines~15 and 16 in the definition of the function \cppLstinline{export_AosBasicTraits\_2()} above. It also saves the class object of the \cppLstinline{Point_2} type (see Line~14), so that later on the functions responsible for the bindings of the various traits models can add additional Python attributes (such as the \cppLstinline{x()}, \cppLstinline{y()}, and \cppLstinline{approximate()} functions) to the class object.

The concept \cppLstinline{AosHorizontalSideTraits\_2} requires the provision of the   \cppLstinline{Compare\_ x\_on\_boundary\_2} functor,\footnote{A functor is a class that has one or more \cppLstinline{operator()} members; thus, it acts as a function.} which has two versions of the overloaded members \cppLstinline{operator()}: one accepts two ends of curves (refer to the manual for the precise encoding of a curve end) and one accepts an end of a curve end and a point. The former compares the $x$-coordinates of the curves at their respective limits or ends and the latter compares the $x$-coordinates of the point and the $x$-coordinates of the curve at its respective limit or end.
%
The concept \cppLstinline{ArrangementIdentifiedHorizontalTraits\_2} indirectly refines the concept \cppLstinline{AosHorizontalSideTraits\_2} mentioned above. In addition to the two versions of the overloaded members \cppLstinline{operator()} of the functor  \cppLstinline{Compare\_x\_on\_boundary\_2} it requires the provision of a third version that accepts two points and compares their $x$ coordinates. Similar to the case described in the previous paragraph, where the class object that handles the bindings for the \cppLstinline{Point_2} type must be available in two separate functions, also here the class object that handles the bindings for the functor \cppLstinline{Compare\_x\_on\_boundary\_2} must be available in both functions \cppLstinline{export_AosHorizontalSideTraits\_2()} and  \cppLstinline{export_ArrangementIdentifiedHorizontalTraits\_2()}. Saving the class object in one function so that it can be reused latter on in the other addresses this necessity.

% -----------------------------------------------------------------------------
\subsection{Return value policy}
\label{sec:code:rvp}
% -----------------------------------------------------------------------------
Python and C++ use fundamentally different ways of managing the memory and lifetime of objects. This can lead to issues when creating bindings for functions that return a non-trivial type. Just by looking at the type, it is unclear whether Python should take charge of the returned value and eventually free its resources, or whether this is handled by the C++ code. When writing code in C++, it is usually considered a good practice to use smart pointers, which exactly describe ownership semantics. Still, even good C++ interfaces use raw references and pointers sometimes. In some cases, in order to assure proper memory management of a return value from a function, explicitly specifying a return value policy is needed. The policy may depend on whether the returned object is a newly created object, a reference to some already existing (internal) object, or a const reference to some existing object. Additionally, in some cases we need to tie the lifetime of the result to the lifetime of the arguments, or the other way around.

For some functions of kernel objects the kind of return value depends on the kernel type; thus, the return value policy needs to be set accordingly. For example, when the kernel type is
\cppLstinline{Exact\_predicates\_inexact\_constructions\_kernel}, the return type of some functions that return a coordinate or a coefficient is a const reference, while when the kernel type is \cppLstinline{Exact\_predicates\_exact\_constructions\_kernel} the same functions return an object by value.

This is why for these functions a \myLstinline{Kernel_return_value_policy} type is being passed as the return value policy. The type is defined differently (via preprocessor macros) based on the kernel being exposed.

% -----------------------------------------------------------------------------
\subsection{Scopes}
\label{sec:code:scopes}
% -----------------------------------------------------------------------------
The Boost.Python mechanism that depends on \cppLstinline{bp::scope} is used to (i) reflect nesting in C++ code that uses \cgal, and (ii) prevent name conflicts between bound entities in distinct modules; see Section~\ref{sec:gen:scopes}. The latter requires a slight intervention on our part as explained next. The code that binds all functions and classes of different modules is nested in distinct C++ forced scopes. Within a (C++) scope we first set the Python namespace to match the name of a module. Then, we bind all functions and classes of \cgal{} packages associated with the module. The following code except shows how it is done for the \iiDandiiiDGeometryKernelPackage{}
module:

\begin{lstlisting}[style=cppNumberedStyle]
{
  SET_SCOPE("Ker")
  export_kernel();
#ifdef CGALPY_KERNEL_INTERSECTION_BINDINGS
  export_kernel_intersections_2();
  export_kernel_intersections_3();
#endif
}
\end{lstlisting}

For every module there exsists a function called \cppLstinline{export_<name>}, where \cppLstinline{<name>} is the module. This function contains the binding code for the essential
functions and classes of the module. Recall that the \iiDandiiiDGeometryKernelPackage{} module is associated with the \myLstinline{KERNEL\_INTERSECTION\_BINDINGS} flag, which determines
whether to generate bindings for intersections; hence, the compile-time conditional calls to the functions \cppLstinline{export_kernel_intersections_2()} and \cppLstinline{export_kernel_intersections_3()}, which contains the binding code for intersections.

%% All functions and classes related to some package were exposed inside
%% a \myLstinline{export_[package]()} function. All such calls are
%% happening in \myLstinline{export_module.cpp}. We wrap each call with a
%% local (c++) scope and call the \myLstinline{SET_SCOPE(module_name)}
%% macro, which causes every exposed function/class inside this local
%% scope to be exposed inside the \myLstinline{module_name} submodule of
%% the CGALPY module. This prevents name clashes of classes with the same
%% name in different packages.

\begin{comment}
% -----------------------------------------------------------------------------
\subsection{Avoiding Code Duplication}
\label{sec:code:templates}
% -----------------------------------------------------------------------------
Some types have the exact same nested types and member functions exposed, e.g., \cppLstinline{Triangulation\_2} and \cppLstinline{Delaunay\_triangulation\_2} class templates each have the method \cppLstinline{dimension()} exposed as the Python attribute \pyLstinline{dimension}. In these cases we use a function template to avoid code duplication. Following the example above, we have the function \cppLstinline{export\_triangulation\_2<Triangulation>()}, where the template parameter \cppLstinline{Triangulation} is substituted with one of the two-dimensional triangulation types
supported; see Table~\ref{tab:tri2:names}. This function template accepts as its single parameter the binding class object, which is constructed separately.

\begin{lstlisting}[style=cppNumberedStyle]
typedef CGAL::Triangulation_2<Kernel>              Triangulation_2;
typedef CGAL::Delaunay_triangulation_2<Kernel>     Delaunay_triangulation_2;

auto c0 = bp::class_<Triangulation_2>("Triangulation_2");
auto c1 = bp::class_<Delaunay_triangulation_2,
                     bp::bases<Triangulation_2>>("Delaunay_triangulation_2");

export_triangulation_2<Triangulation_2>(c0);
export_triangulation_2<Delaunay_triangulation_2>(c1);
\end{lstlisting}
\end{comment}

% -----------------------------------------------------------------------------
\subsection{Substitution Failure Is Not an Error}
\label{sec:code:sfinae}
% -----------------------------------------------------------------------------
Substitution failure is not an error (SFINAE) refers to a situation in C++ where an invalid substitution of template parameters is not in itself an error. In particular, it applies during a process called \emph{overload resolution}.  In order to compile a function call, the compiler creates a set of candidate functions. This set includes all function implementations the name of which match the call and that can be accessed by the caller. Then the compiler selects the most viable one according to some criteria. When substituting the explicitly specified or deduced type for the template parameter fails, the specialization is discarded from the overload set instead of causing a compile error.\footnote{For more information see, e.g., \url{https://en.cppreference.com/w/cpp/language/sfinae}.}

The function \cppLstinline{CGAL::do_intersect(o1, o2)} is overloaded with several implementations that handle different combination of types of arguments; however, not every combination of two types is implemented. For a complete list of valid combinations refer to the reference manual at
\url{https://doc.cgal.org/latest/Kernel_23/group__do__intersect__linear__grp.html}. We use SFINAE to create bindings for all supported combinations and avoid getting compilation errors for unsupported combinations as follows. Consider two types \cppLstinline{Type1} and
\cppLstinline{Type2}, and assume that the overloaded version \cppLstinline{CGAL::do_intersect(Type1 o1, Type2 o2)} is not implemented but all other three are implemented.

\begin{lstlisting}[style=cppNumberedStyle,basicstyle=\scriptsize]
template <typename T1, typename T2>
void bind_do_intersect_2T(decltype(CGAL::do_intersect<Kernel>(T1(), T2()))) {
  bp.def<bool(const T1&, const T2&)>("do_intersect", &CGAL::do_intersect<Kernel>);
}

template <typename, typename>
void bind_do_intersect_2T(...) {}

template <typename T1>
void bind_do_intersect_1T() {
  bind_do_intersect_2T<T1, Type1>(true);
  bind_do_intersect_2T<T1, Type2>(true);
}

void bind_do_intersect() {
  bind_do_intersect_1T<Type1>();
  bind_do_intersect_1T<Type2>();
}
\end{lstlisting}

First the function \cppLstinline{bind_do_intersect()} calls the instance of \cppLstinline{bind_do_intersect_1T<T1>()}, where \cppLstinline{T1} is substituted with \cppLstinline{Type1}.  Then, the latter makes two calls and passes a \cppLstinline{bool} argument of value \cppLstinline{true}. (This value can be \cppLstinline{false} just as well.) First it calls the instance of \cppLstinline{bind_do_intersect_2T<Type1, T2>(true)}, where \cppLstinline{T2} is also substituted with \cppLstinline{Type1}; then, it calls the instance of \cppLstinline{bind_do_intersect_2T<Type1, T2>(true)}, where \cppLstinline{T2} is substituted with \cppLstinline{Type2}. The function \cppLstinline{bind_do_intersect_2T()} accepts a parameter of type \cppLstinline{bool} and is overloaded with two implementations. The second uses \emph{variadic arguments} and serves as a fall-through when the evaluation of the first implementation fails during the resolution process.\footnote{For more information on \emph{variadic arguments} see \url{https://en.cppreference.com/w/cpp/language/variadic_arguments}.} The parameter type of the first implementation is defined as the return type of \cppLstinline{CGAL::do_intersect<Kernel>(T1(), T2())}. When \cppLstinline{T1} and \cppLstinline{T2} are substituted with \cppLstinline{Type1} and \cppLstinline{Type2}, respectively, the evaluation fails, and the blank implementation is selected and invoked.  In all other three cases the evaluation succeeds. As variadic parameters have the lowest rank for the purpose of overload resolution, the first implementation will always be selected for valid combinations. Here, the return type of \cppLstinline{CGAL::do_intersect<Kernel>(T1(), T2())} is always resolved at compile time (to \cppLstinline{bool}) when the function is defined.

The situation presented by the function template \cppLstinline{minkowski_sum_2()} that applies the \emph{convex-decomposition approach} (see Section~\ref{ssec:modules:ms2}) is a bit more complicated. Consider for example the function shown at Line~1 of Listing~\ref{lst:msSignatures}. We need to generate bindings for instances of this function only for valid combinations of \cppLstinline{PolygonType1}, \cppLstinline{PolygonType2}, and \cppLstinline{PolygonConvexDecomposition_2} types. In the code below we use the respective short names \cppLstinline{PT1}, \cppLstinline{PT2}, and \cppLstinline{PP} instead. A valid combination exists only if the latter is a polygon-partition unary function that can be applied to a polygon of type \cppLstinline{PolygonType1} and to a polygon of type \cppLstinline{PolygonType2}. In particular, if \cppLstinline{pp} is a polygon-partition function, the calls \cppLstinline{pp.(pgn1, it)} and \cppLstinline{pp.(pgn2, it)} must be valid, where \cppLstinline{pgn1} is the input polygon of type \cppLstinline{PolygonType1}, \cppLstinline{pgn2} is the input polygon of type \cppLstinline{PolygonType2}, and \cppLstinline{it} is an output iterator of a container of the resulting polygons. Dereferencing the iterator must yield a type that can represent a simple polygon. If \cppLstinline{PolygonType1} is identical to \cppLstinline{Polygon_2}, then \cppLstinline{std::list<PolygonType1>::iterator} can serve as an output iterator; otherwise, \cppLstinline{PolygonType1} must represent a polygon with holes, and in this case \cppLstinline{std::list<PolygonType1::Polygon_2>::iterator} can serve as an output iterator. We use SFIAE for the second time to make this distinction.
%
First, we introduce a class template called \cppLstinline{target}; see Line~6 in the Listing below; it is parameterized with a named parameter \cppLstinline{T} and an unnamed parameter that defaults to \cppLstinline{void}. It delegates the type \cppLstinline{std::list<T>::iterator}. We also introduce a partial  specialization (see Line 9) that delegates the type \cppLstinline{std::list<T::Polygon_2>::iterator}. The second parameter is implicitly substituted with a type provided by a utility template class called \cppLstinline{if_}; see Lines~1--4 in the Listing below. If the type \cppLstinline{T::Polygon_2} is not defined, the instance of the specialization is discarded from the overloaded resolution set. Otherwise, it is selected over the default implementation, as it has a higher rank.
%
We need to capture valid combinations of three types. Recall that the code listed in the previous section captures valid combinations of two types. Following the idea presented in the previous section, we introduce three template functions, where \cppLstinline{bind\_mink\_sum\_decomp\_one\_strategy_3T()} is the last one called; see Line~14 for the default implementation and Line~17 for the partially specialized implementation in the listing below. The polygon-partition function returns an output iterator that points to next to the last element of the container. Passing an argument the type of which matches the type of the returned value of the polygon-partition function to the function \cppLstinline{bind\_mink\_sum\_decomp\_one\_strategy_3T()} is complicated, as this type is not a simple type (e.g., \cppLstinline{bool}) any more. Instead of using the return type, we introduce two unnamed template parameters with default values, that evaluate to the return type of the polygon-partition function when applied to a polygon of type \cppLstinline{PolygonType1} or type \cppLstinline{PolygonType1}, respectively; see Lines~15 and 16 in the listing below. Observe that the outcome of the evaluation is irrelevant. The only thing we are concerned about is weather the evaluation succeeds---it succeeds only if the matching function is defined.

\begin{lstlisting}[style=cppNumberedStyle,basicstyle=\scriptsize]
template <bool Condition, typename Then, typename Else = void>
struct if_ { typedef Then type; };
template <typename Then, typename Else>
struct if_<false, Then, Else > { typedef Else type; };

template <typename T, typename = void> struct target {
  typedef typename std::list<T>::iterator type;
};
template <typename T>
struct target<T, typename if_<false, typename T::Polygon_2>::type> {
  typedef typename std::list<typename T::Polygon_2>::iterator type;
};

template <typename T1, typename T2, typename T3>
void bind_mink_sum_decomp_one_strategy_3T(...) {}

template <typename PT1, typename PT2, typename PP,
          typename = decltype(PP()(PT1(), typename target<PT1>::type())),
          typename = decltype(PP()(PT2(), typename target<PT2>::type()))>
void bind_mink_sum_decomp_one_strategy_3T(bool) {
  bp::def<Polygon_with_holes_2(const PT1&, const PT2&, const PP&)>
    ("minkowski_sum_2", &CGAL::minkowski_sum_2<Kernel, Point_2_container, PP>);
}
\end{lstlisting}

%% templated with some supported type
%% {T1}. \myLstinline{bind_do_intersect_1T()}, templated with
%% \myLstinline{T1}, calls \myLstinline{bind_do_intersect_2T} templated
%% with \myLstinline{T1} and some supported type \myLstinline{T2}, and a
%% \myLstinline{bool} argument.  The declaration of
%% \myLstinline{bind_do_intersect_2T}, templated with both
%% \myLstinline{T1} and \myLstinline{T2} is only valid when templated
%% with a combination of types for which an overloaded implementation of
%% \myLstinline{CGAL::do_intersect} exists. Otherwise, due to SFINAE, the
%% blank implementation template of \myLstinline{bind_do_intersect_2T}
%% expecting (...) as an argument will be invoked (it would normally not
%% be invoked as variadic parameters have the lowest rank for the purpose
%% of overload resolution).

We use SFINAE to efficiently generate bindings of other types as described in the following sections.

% -----------------------------------------------------------------------------
\subsection{Arrangement Extension}
\label{sec:code:ae}
% -----------------------------------------------------------------------------
In many applications it is necessary to extend the type that represents certain \cgal{} data structures with some types that represent new properties. In \cgal{} it is more convenient, and thus more common, extending the types that represent the topological features of the data structures (rather then the types that represent the geometric features), for example, the DCEL of the arrangement or the triangulation data structure of a triangulation. When developing code in pure C++, there is no restriction on the property types used to extend the data structure. However, when the data structure interface is wrapped with Python bindings, the property type must be known when the bindings are compiled. Nevertheless, we support extensions with generic Python objects. If a certain feature must be extended with a certain property, the type of which is defined in the C++ code, it is always possible to use an external property map, where the properties are indexed by a Python object, such as an integer or a string. When generating the bindings the user must set the flag that enables the extension (unless extension is the default setting). Attaching indices
to and extracting an index from a feature is done via the \cppLstinline{set_data()} and \cppLstinline{data()} exposed functions. Each cell type, namely, \cppLstinline{Vertex}, \cppLstinline{Halfedge}, or \cppLstinline{Face}, can be extended independently (see Table~\ref{tab:aos}), which implies that there are eight plausible combinations that the user can select form to form the final DCEL type of the arrangement.

In the following we assume that the face type of the \cppLstinline{Arrangement_2} instance is extended. The Python code excerpt below constructs two arrangements. For each arrangement the properties of the unbounded face and the other face are initialized with a Python integer objects 0 and 1, respectively. The code computes their overlay (see Figure~\ref{fig:faceExt}) and updates the face property to indicate the number of overlapping bounded faces.

\begin{lstlisting}[style=pythonFramedStyle,basicstyle=\scriptsize,label={lst:arrExtend},caption={Constructing two extended arrangements and initializing the extended data.}]
CGALPY = importlib.import_module(lib)
Aos2 = CGALPY.Aos2
Arrangement_2 = Aos2.Arrangement_2
Point_2 = Arrangement_2.Geometry_traits_2.Point_2
Curve_2 = Arrangement_2.Geometry_traits_2.Curve_2

arr1 = Arrangement_2()
c1 = Curve_2(Point_2(0, 0), Point_2(2, 0))
c2 = Curve_2(Point_2(2, 0), Point_2(2, 2))
c3 = Curve_2(Point_2(2, 2), Point_2(0, 2))
c4 = Curve_2(Point_2(0, 2), Point_2(0, 0))
Aos2.insert(arr1, [c1, c2, c3, c4])

arr2 = Arrangement_2()
c1 = Curve_2(Point_2(1, 1), Point_2(3, 1))
c2 = Curve_2(Point_2(3, 1), Point_2(3, 3))
c3 = Curve_2(Point_2(3, 3), Point_2(1, 3))distance=-1pt
c4 = Curve_2(Point_2(1, 3), Point_2(1, 1))
Aos2.insert(arr2, [c1, c2, c3, c4])

for arr in [arr1, arr2]:
  for f in arr.faces():
    f.set_data(0) if f.is_unbounded() else f.set_data(1)
\end{lstlisting}

\begin{figure}[!htp]
  \captionsetup[subfigure]{justification=centering}
  \centering%
  \subfloat[]{\label{fig:faceExt}
    \begin{tikzpicture}
      \fill[polyALightColor] (0,0) rectangle (2,2);
      \fill[polyBLightColor,opacity=0.5] (1,1) rectangle (3,3);
      \draw[thick,polyADarkColor] (0,0) rectangle (2,2);
      \draw[thick,polyBDarkColor] (1,1) rectangle (3,3);
      \node at (2.5,0.5) {[0]};
      \node at (0.5,0.5) {[1]};
      \node at (2.5,2.5) {[1]};
      \node at (1.5,1.5) {[2]};
      \node[point] at (0,0) {};
      \node[point] at (2,0) {};
      \node[point] at (0,2) {};
      \node[point] at (2,2) {};
      \node[point] at (1,1) {};
      \node[point] at (3,1) {};
      \node[point] at (1,3) {};
      \node[point] at (3,3) {};
      \node[point] at (2,1) {};
      \node[point] at (1,2) {};
    \end{tikzpicture}}\quad\quad
  \subfloat[]{\label{fig:vertexExt}
    \begin{tikzpicture}[scale=0.2]
      \draw[thick,polyADarkColor] (3,4) circle (5);
      \draw[thick,polyADarkColor] (-3,-4) circle (5);
      \draw[thick,polyBDarkColor] (4,-3) circle (5);
      \draw[thick,polyBDarkColor] (-4,3) circle (5);
      \node[point=red!50!white,label={[label distance=-1pt]180:4}] at (-8, -4) {}; % 4
      \node[point=green,label={[label distance=-1pt]180:6}] at (-7, -1) {};        % 6
      \node[point=red!50!white,label={[label distance=-1pt]180:2}] at (-9, 3) {};  % 2
      \node[point=red!50!white,label={[label distance=-1pt]180:4}] at (-2, 4) {};  % 4
      \node[point=green,label={[label distance=-1pt]90:6}] at (-1, 7) {};          % 6
      \node[point=red!50!white,label={[label distance=-1pt]180:2}] at (-1, -3) {}; % 2
      \node[point=blue,label={[label distance=-1pt]0:12}] at (0, 0) {};            % 12
      \node[point=green,label={[label distance=-1pt]-90:6}] at (1, -7) {};         % 6
      \node[point=red!50!white,label={[label distance=-1pt]0:2}] at (1, 3) {};     % 2
      \node[point=red!50!white,label={[label distance=-1pt]0:4}] at (2, -4) {};    % 4
      \node[point=green,label={[label distance=-1pt]0:6}] at (7, 1) {};            % 6
      \node[point=red!50!white,label={[label distance=-1pt]0:4}] at (8, 4) {};     % 4
      \node[point=red!50!white,label={[label distance=-1pt]0:2}] at (9, -3) {};    % 2
    \end{tikzpicture}}
  \caption{%
    (\subref{fig:faceExt}) The arrangement faces are extended with integers that indicate the number of overlapping bounded faces, shown in brackets.
    (\subref{fig:vertexExt}) The arrangement vertices are extended with integers that indicate the weighted degree of the vertices.}
  \label{fig:extend}
\end{figure}

Extending the arrangement features is useful for many applications, but essential for applications that compute the overlay of arrangements as discussed in the next Section.

\begin{comment}
- Extending the 2d and 3d vertex and face triangulation types => change the title to be more general in the full article
\end{comment}

% -----------------------------------------------------------------------------
\subsection{Arrangement Overlay Traits}
\label{sec:code:overlay}
% -----------------------------------------------------------------------------
The map overlay of two arrangements $\calA_1$ and $\calA_2$, conveniently referred to as the \textcolor{red}{red} and \textcolor{blue}{blue} arrangements, is a third arrangement $\calA$, such that there is a cell $c$ in $\calA$ iff there are cells $c_1$ and $c_2$ in $\calA_1$ and $\calA_2$, respectively, such that $c$ is a maximal connected component of $c_1 \cap c_2$. Arrangements are subdivisions of some ambient space (see Section~\ref{ssec:modules:aos2}) and computing the overlay of two arrangements is useful for many applications, including the application of Boolean operations; see Section~\ref{ssec:modules:bso2}. Indeed, the \iiDArrangementsPackage{} package supports an operation that computes the map overlay of two arrangements in form of the free function template \cppLstinline{CGAL::overlay()}. In particular, the call \cppLstinline{CGAL::overlay(arr_r, arr_b, arr)} computes the overlay of the arrangements \cppLstinline{arr_r} and \cppLstinline{arr_b} and stores the result in the arrangement \cppLstinline{arr}. All three types must be instances of the \cppLstinline{Arrangement_2<Traits, Dcel>} class template. When the function is used in pure C++ code, the geometry-traits classes that substitute the \cppLstinline{Traits} template parameters of the input arrangements must be convertible to the geometry-traits class of the resulting arrangement. We only wrap instances of the overlay function with Python bindings, where all three arrangements use the exact same geometry-traits classes, and thus the types of the three arrangements of the exposed \cppLstinline{overlay()} function are identical.

The \cppLstinline{overlay()} function template is overloaded with a variant that accepts four arguments. The last argument, referred to as the overlay traits, enables the update of properties of the output arrangement from the properties of the input arrangements. Using the overlay traits is typically accompanied with the extension of one or more types of the arrangement cells that hold the properties. The overlay traits must model the \cppLstinline{OverlayTraits} concept, which requires the provision of ten functions that handle all possible cases as listed below. Let $v_r$, $e_r$, and $f_r$ denote input \textcolor{red}{red} cells, i.e.,~a vertex, an edge, and a face, respectively, $v_b$, $e_b$, and $f_b$ denote input \textcolor{blue}{blue} features, and $v$, $e$, and $f$ denote output cells.
\begin{enumerate}
\item A new vertex $v$ is induced by coinciding vertices $v_r$ and $v_b$.
\item A new vertex $v$ is induced by a vertex $v_r$ that lies on an edge $e_b$.
\item An analogous case of a vertex $v_b$ that lies on an edge $e_r$.
\item A new vertex $v$ is induced by a vertex $v_r$ that is contained in a face $f_b$.
\item An analogous case of a vertex $v_b$ contained in a face $f_r$.
\item A new vertex $v$ is induced by the intersection of two edges $e_r$ and $e_b$.
\item A new edge $e$ is induced by the (possibly partial) overlap of two edges $e_r$ and $e_b$.
\item A new edge $e$ is induced by the an edge $e_r$ that is contained in a face $f_b$.
\item An analogous case of an edge $e_b$ contained in a face $f_r$.
\item A new face $f$ is induced by the overlap of two faces $f_r$ and $f_b$.
\end{enumerate}
Evidently a custom C++ overlay traits cannot be defined in Python; more precisely, the traits type must be known when the bindings are compiled. We introduce and expose with Python bindings two models of the \cppLstinline{OverlayTraits} concept. The most general among the two, called \cppLstinline{Arr_overlay_traits}, defines ten functions that correspond to the list above, each accepting three arguments, namely, two objects that represent input cells, and one object that represents an output cell. By default these functions are idle. A user (that is, a Python programmer) can override any subset of these functions. One constructor of this model accepts all the ten functions at once. Another constructor accepts a single function, which corresponds to function (10) in the list above. In the following we assume that the face type of the \cppLstinline{Arrangement_2} instance is extended. The code excerpt below computes the overlay of the two arrangements constructed by the code in Listing~\ref{lst:arrExtend} using the overlay traits. Recall that the properties that comprise the extension are represented by Python objects to start with; see Section~\ref{sec:code:ae} and that the property of each face in the input arrangement is initialized with an integer Python object. The property of each face of the resulting arrangement is then set with the sum of the integers associated with its overlapping faces during the overlay computation process.

\begin{lstlisting}[style=pythonFramedStyle,basicstyle=\scriptsize]
traits = Aos2.Arr_overlay_traits(lambda f1,f2,f: f.set_data(f1.data()+f2.data()))
Aos2.overlay(arr1, arr2, result, traits)
\end{lstlisting}

Observe that while the computation of the overlay is carried out by the compiled C++ code; however, the summation is carried out by a Python \pyLstinline{lambda} function, which accepts as the third argument the face to be updated. By default, arguments are copied into new Python objects. We change this behaviour and pass the third object by value to ensure that the update is persistent.

In many cases, such as in the example above, the property of the new cell depends solely on the properties of the overlapping two cells. To this end we introduce (and expose) a second model of the \cppLstinline{OverlayTraits} concept called \cppLstinline{Arr_overlay_function_traits}. This model also defines ten functions that correspond to the list above. However, here each function accepts the two Python objects that extends the overlapping cells as input arguments and returns the Python object that extends the resulting cell. The code excerpt below has the same affect  the code above has, but is more compact.

\begin{lstlisting}[style=pythonFramedStyle,basicstyle=\scriptsize]
traits = Aos2.Arr_overlay_function_traits(lambda x, y: x+y)
os2.overlay(arr1, arr2, result, traits)
\end{lstlisting}

Each one of the ten functions supported by the \cppLstinline{Arr_overlay_function_traits} model returns a Python object that is passed back to the compiled C++ code. Then, the object is stored with the new cell. Thus, there is no need to pass arguments by reference (even thought it could be more efficient in certain cases); however, the implementation of this model presents a new coding challenge.

Assume that only the vertex type of the \cppLstinline{Arrangement_2} instance is extended. The Python code excerpt below computes the overlay of two arrangements and updates the property of each vertex of the resulting arrangement to indicate the weighted degree of the vertex, where blue incident edges weigh twice as much as red incident edges; see Figure~\ref{fig:vertexExt}. We set all the six functions that update output vertices in the \cppLstinline{Arr_overlay_function_traits} object using dedicated setters. The name of a setter matches the pattern \lstinline[language={C++},mathescape=true,columns=fixed]{set_$c_rc_b$_$c$()}, where each of $c_r$, $c_b$, and $c$ can be substituted with \cppLstinline{v}, \cppLstinline{e}, or \cppLstinline{f}. $c_r$ and $c_b$ determine the input red and blue cell types, respectively, and $c$ determines the output cell type.

\begin{lstlisting}[style=pythonFramedStyle,basicstyle=\scriptsize]
Aos2.insert(arr1,[Curve(Circle(Point(3,4),25)),Curve(Circle(Point(-3,-4),25))])
Aos2.insert(arr2, [Curve(Circle(Point(3,4),25)), Curve(Circle(Point(3,4),25))])
for arr in [arr1, arr2]:
  for v in arr.vertices():
    v.set_data(v.degree())
traits = Aos2.Arr_overlay_function_traits()
traits.set_vv_v(lambda x, y: 2*x+y)
traits.set_ve_v(lambda x, y: 2*x+2)
traits.set_vf_v(lambda x, y: 2*x)
traits.set_ev_v(lambda x, y: 4+y)
traits.set_fv_v(lambda x, y: y)
traits.set_ee_v(lambda x, y: 6)
Aos2.overlay(arr1, arr2, result, traits)
\end{lstlisting}

Consider a generic function called \cppLstinline{apply()} that accepts two objects that represent input cells, namely \cppLstinline{r} and \cppLstinline{b}, one object that represents an output cell, namely \cppLstinline{o}, and a function called \cppLstinline{f}. The objective of \cppLstinline{apply()} is to apply \cppLstinline{f} to the properties of the cells \cppLstinline{r} and \cppLstinline{b} and store the return value as the property of the cell \cppLstinline{o}. Recall that a property of a cell \cppLstinline{c} is obtained and set via the calls \cppLstinline{c.data()} and \cppLstinline{c.set_data()}, respectively. If the type of the output cell is not extended, however, \cppLstinline{apply()} should become idle. Similarly, if the type of an input cell is not extended, the \pyLstinline{None} python object (represented by \cppLstinline{bp::object}) should be passed to \cppLstinline{f} instead. We use SFINAE yet again, several times, to address the above. First, we introduce an overload of \cppLstinline{apply()} that is idle and serves as a fall-through; see Line~9 in the listing below. When a call to \cppLstinline{apply()} is made. The idle overload is selected by the overload resolution process if the output cell does not have a member called \cppLstinline{set_data()}. If the output cell does have this member the other version is selected and calls \cppLstinline{data_a(r)} and \cppLstinline{data_a(b)} to obtain the properties of the red and blue input cells, respectively. Second, we introduce an overload of \cppLstinline{data_a()} that serves as a fall-through; see Lines 2. It does nothing but return the \pyLstinline{None} Python object. It is selected if the corresponding cell does not have a member called \cppLstinline{data()}. If the cell does have this member the other version is selected; it returns the property of the cell.

\begin{lstlisting}[style=cppFramedStyle,basicstyle=\scriptsize]
// Fall-through; A::data() does not exist
template <typename A> bp::object data_a(...) { return bp::object(); }

// A::data() exists
template <typename A, typename = decltype(std::declval<A>().data())>
const bp::object& data_a(const A* a) { return a->data(); }

// Fall-through; O::set_data() does not exist
template <typename R, typename B, typename O, typename F> void apply(...) {}

// O::set_data() does exist
template <typename R, typename B, typename O, typename F,
          typename = decltype(std::declval<O>().set_data(std::declval<typename O::Data>()))>
void apply(const R* r, const B* b, O* o, F f) {
  o->set_data(f(data_a<R>(r), data_a<B>(b)));
}
\end{lstlisting}

%% (the extended arrangement data type) Those 10 functions are saved as
%% class members and are then called in the respective 10 [overloaded]
%% versions of \myLstinline{create_vertex}, \myLstinline{create_edge},
%% \myLstinline{create_face object}) (the matching is done by the order
%% of the passed functions and the order of the [overloaded] versions in
%% the official concept reference).

%% This way the user can still define the way the result face/edge/vertex
%% data will be calculated.

\begin{comment}
% -----------------------------------------------------------------------------
\subsection{kd tree}
\label{sec:code:kdtree}
% -----------------------------------------------------------------------------
The k-d tree dimension and the dimension of related classes is templated and can thus only be decided at compile time via the cmake variable \myLstinline{CGALPY_SPATIAL_SEARCHING_DIMENSION}.

In order to assert the dimension the bindings were compiled with, a free function \myLstinline{get_spatial_searching_dimension()} was exposed that returns the dimension as an \myLstinline{int}.

% -----------------------------------------------------------------------------
\subsection{\myLstinline{Python_distance}}
\label{sec:code:distance}
% -----------------------------------------------------------------------------
Similarly to overlay traits, a user may want to define a custom distance class to use with search queries. Here too, as defining a new C++ class is not possible from python, we provide a similar solution: A \myLstinline{Python_distance} class is exposed, which is a model of
\myLstinline{GeneralDistance} concept. \url{https://doc.cgal.org/5.2/Spatial_searching/classGeneralDistance.html} It expects 5 functions in its constructor, and those are saved as class members and are used for the respective 5 operations [the optional are not included] of the concept.

Making some things more pythonic:
\end{comment}

% -----------------------------------------------------------------------------
\subsection{Iterators and circulators}
\label{sec:code:iters}
% -----------------------------------------------------------------------------
For binding iterators we used the \myLstinline{bind_iterator} templated function and its variations, in order to allow pythonic iteration. Circulators too are artificially converted to iterators supporting \myLstinline{__next__} and throwing a stop iteration when reaching the first item again, by being wrapped in a templated \myLstinline{Iterator_from_circulator} object which is the one being exposed to python.

% -----------------------------------------------------------------------------
\subsection{Functions populating an output iterator}
\label{sec:code:fnc-out}
% -----------------------------------------------------------------------------
Instead, a wrapper that expects a python list and populates it with the results is being exposed.

% -----------------------------------------------------------------------------
\subsection{Functions expecting begin and end iterators}
\label{sec:code:fnc-range}
% -----------------------------------------------------------------------------
Instead, a wrapper that expects a python list as an argument and calls the function with the list's begin and end iterators is exposed.

% -----------------------------------------------------------------------------
\subsection{Type Annotation}
\label{sec:code:typeAnnotation}
% -----------------------------------------------------------------------------
\cgal{} is a large library to start with. It is divided into approximately~150 packages with countless function and class templates. The vast number of function and type instances that can be defined using this templates can be overwhelming. The development of non critical code that is based on \cgal{} can be expedited using Python and \cgal{} Python bindings.  Using code completion (which is based on type annotations) offered by several Python IDEs can accelerate the process much further. Type annotation in Python refers to annotations of Python functions, arguments, and variables in a way that can be used by various tools. It is an optional feature of Python that has been introduced in Version~3.5 and enhanced in successive version. With this feature implemented a type checker, a tool separate from the Python interpreter, can be used to statically analyze code base to find bugs during early development stages and spelunk code of large projects. Naturally, it can be used by IDEs to interactively provide hints and suggestions to programmers and type check the code as it is being developed. Annotations can and should be placed near the code if possible. Observe that the Python interpreter ignores these annotations. Annotations can also be placed in external files (that have the extension \lstinline{pyi}), referred to as stubs, for cases where the code base is inaccessible, such as the case with our bindings. Stub files contain the signatures of the annotated functions with the body of the functions discarded. They also list the attributes of annotated types. In our case they contain the annotated interface of the \cgal{} bindings. Stub files have the same syntax as regular Python modules. There is one feature of the typing module that is different in stub files, namely, the \pyLstinline{@overload} decorator. This decorator allows describing functions that support multiple different combinations of argument types. Consider, for example, the overloaded free C++ functions \cppLstinline{overlay()}; see Section~\ref{sec:code:overlay}. The annotation for the corresponding Python function \pyLstinline{overlay()} in the stub file follows.
\begin{lstlisting}[style=pythonFramedStyle,basicstyle=\scriptsize]
Aos2 = CGALPY.Aos2
Arr = Aos2.Arrangement_2
@overload
def Aos2.overlay(r: Arr, b: Arr, o: Arr) -> None: pass
@overload
def Aos2.overlay(r: Arr, b: Arr, o: Arr, t: Aos2.Arr_overlay_traits) -> None: pass
@overload
def Aos2.overlay(r: Arr, b: Arr, o: Arr, t: Aos2.Arr_overlay_function_traits) -> None: pass
\end{lstlisting}

Annotating a Python type that wraps a C++ type, which is a model of some concepts is guided by the concepts. Similar to the C++ code of the bindings, we have introduced a tight coupling between concepts and Python type annotations; see Section~\ref{sec:code:concepts}. In particular, we have introduced a Python annotation class for each concept that annotates all the requirements of the concept. The refinement relations between concepts are not reflected using inheritance between concept annotation classes. (This would invalidate the annotation code, as a concept \cppLstinline{A} may refine concepts \cppLstinline{B} and \cppLstinline{C} and both, \cppLstinline{B} and \cppLstinline{C}, may refine concept \cppLstinline{D}.) Instead, every model annotation class that annotates a certain type, say \cppLstinline{T}, inherits from all the concept annotation classes that correspond to the concepts that model \cppLstinline{T}. For example, assume that bindings are generated for the \cppLstinline{Arrangement_2<GeometryGtaits_2, Dcel>} type instantiated with a the geometry traits type that is an instance of the \cppLstinline{Arr_segment_traits_2} class template. This type models several concepts; , e.g., Figure~\ref{fig:atc:central}. Partial Python type annotation for the type \pyLstinline{Arr_segment_traits_2} follows.
\begin{lstlisting}[style=pythonFramedStyle,basicstyle=\scriptsize]
class Arr_segment_traits_2(AosBasicTraits_2._AosBasicTraits_2,
                           AosXMonotoneTraits_2._AosXMonotoneTraits_2,
                           AosTraits_2._AosTraits_2,
                           AosLandmarkTraits_2._AosLandmarkTraits_2,
                           AosVerticalSideTraits_2._AosVerticalSideTraits_2,
                           AosVerticalSideTraits_2._AosHorizontalSideTraits_2) :
  class Point_2(AosBasicTraits_2._AosBasicTraits_2.Point_2):
    def x(self) -> Ker.FT: ...
    def y(self) -> Ker.FT: ...
\end{lstlisting}
The type annotation for the concept \cppLstinline{AosBasicTraits_2} contains the annotations for the default constructor and copy constructor of the type \cppLstinline{Point_2} defined by every model of the concept \cppLstinline{AosBasicTraits_2}.
\begin{lstlisting}[style=pythonFramedStyle,basicstyle=\scriptsize]
class _AosBasicTraits_2():
  class Point_2():
    def __init__(): ...
  class X_monotone_curve_2():
    def __init__(): ...
  class Equal_2():
    @overload
    def __call__(self, p: _AosBasicTraits_2.Point_2, q: _AosBasicTraits_2.Point_2) -> bool: ...
    @overload
    def __call__(self, p: _AosBasicTraits_2.X_monotone_curve_2, q: _AosBasicTraits_2.X_monotone_curve_2) -> bool: ...
  def equal_2_object(self) -> _AosBasicTraits_2.Equal_2: ...
\end{lstlisting}
